%----------------------------------------
\begin{frame}
  \frametitle{Part 1}

{\Large\bf Motivation: why concurrency defeats testing}

\end{frame}

%----------------------------------------
\begin{frame}[fragile]
  \frametitle{Six lines of obviously correct code} \small

Two instruments share a $10\,$A power bus. Each draws $6\,$A, and each asks the
power manager before switching on.

\begin{lstlisting}[language=promela]
byte avail   = 10;             /* SHARED: amps left on the bus  */
byte powered = 0;              /* SHARED: instruments now on    */

inline request_power(amps) {
  if
  :: avail >= amps ->          /* CHECK: is there enough?       */
       avail = avail - amps;   /* ACT:   reserve it             */
       powered++               /*        switch on              */
  :: else -> skip              /* denied - stay off             */
  fi
}

active proctype InstrumentA() { request_power(6) }   /* thread 1 */
active proctype InstrumentB() { request_power(6) }   /* thread 2 */
\end{lstlisting}

\end{frame}

%----------------------------------------
\begin{frame}
  \frametitle{The interleaving that flies} \small

\begin{center}\scriptsize
\begin{tabular}{c l l r r}
\toprule
{\bf step} & {\bf Instrument A} & {\bf Instrument B} & {\bf avail} & {\bf powered}\\
\midrule
0 & \multicolumn{2}{l}{\em (initial)}                       & 10 & 0\\
1 & \tt avail >= 6 ? yes  &                                 & 10 & 0\\
\rowcolor{red!8}
2 &                       & \tt avail >= 6 ? yes            & 10 & 0\\
3 & \tt avail = 10 - 6    &                                 &  4 & 0\\
4 & \tt powered++         &                                 &  4 & 1\\
5 &                       & \tt avail = 4 - 6               & \bf 254 & 1\\
6 &                       & \tt powered++                   & \bf 254 & \bf 2\\
\bottomrule
\end{tabular}
\end{center}

\pause
\vspace{0.3cm}

\begin{block}{Two failures for the price of one}\small
\begin{itemize}
\item {\bf $12\,$A drawn from a $10\,$A bus.} Brownout, undervoltage trip, or a
      reset --- on hardware, not in the model.
\item {\bf \ipromela{avail} underflowed} $4-6 = -2$, and as a \ipromela{byte} that is
      $254$. The power manager now believes it has $254\,$A to give away.
      {\em Every subsequent request succeeds.}
\end{itemize}
\end{block}

\pause

The second failure is the dangerous one: the system does not crash here.
It quietly stops enforcing the budget, and fails somewhere else, much later.

\end{frame}

%----------------------------------------
\begin{frame}
  \frametitle{Why testing does not find this} \small

The window is {\bf two instructions wide} --- it needs a preemption landing
exactly between the check and the act.

\pause
\vspace{0.4cm}

\begin{block}{}\small
\centering
The defect is in the {\bf interleaving}. Reviewing interleavings is\\
far harder than reviewing sequential behavior.
\end{block}

\end{frame}

%----------------------------------------
\begin{frame}
  \frametitle{The number of interleavings explodes} \small

$n$ tasks, each with $m$ atomic steps, gives

\[
\frac{(n \cdot m)!}{(m!)^n}
\]

distinct interleavings.

\vspace{0.15cm}

\begin{center}
\begin{tabular}{r r r}
\toprule
{\bf tasks} & {\bf steps each} & {\bf interleavings}\\
\midrule
2 &  2 & 6 \\
3 &  3 & 1\,680 \\
4 &  4 & 63\,063\,000 \\
5 &  5 & $\approx 6.2 \times 10^{14}$ \\
\rowcolor{red!8}
6 &  6 & $\approx 2.7 \times 10^{24}$ \\
\bottomrule
\end{tabular}
\end{center}

\pause
\vspace{0.1cm}

\begin{columns}[c]
\begin{column}{0.36\textwidth}
\centering
\starfield[0.74]
\end{column}
\begin{column}{0.60\textwidth}\small
Six tasks of six steps have more interleavings than there are
{\bf stars in the observable universe} ($\approx 2 \times 10^{23}$).

\vspace{0.25cm}
\pause
This is the famous {\bf state explosion problem}, and it is the
central obstacle for model checking. We return to it later.
\end{column}
\end{columns}

\end{frame}

% ===== HIDDEN 2026-09-09 at Klaus's request: Mars Pathfinder, 1997 =====
%%----------------------------------------
%\begin{frame}
%  \frametitle{Mars Pathfinder, 1997} \small
%
%Repeated total system resets on the surface of Mars.
%
%\vspace{0.25cm}
%
%\begin{block}{Priority inversion}\small
%A {\bf low}-priority meteorological task held a mutex on the information bus.
%A {\bf high}-priority bus-management task blocked on that mutex.
%A {\bf medium}-priority communications task then preempted the low-priority one.
%The mutex was never released, the high-priority task missed its deadline, and
%the watchdog reset the spacecraft.
%\end{block}
%
%\pause
%\vspace{0.25cm}
%
%{\bf Why the medium task matters.} Without it the blocking is bounded by the
%{\em length of the critical section}. The medium task starves the holder, so the
%blocking is bounded only by how long an {\em unrelated} task chooses to run.
%
%\pause
%\vspace{0.25cm}
%
%{\bf Why it is hard to find.} The medium task {\em never touches the mutex}.
%It appears in no analysis of that lock, and reviewing the two tasks that share
%it reveals nothing. The fault is not in any component --- it is in the global
%interleaving of three.
%
%\pause
%\vspace{0.25cm}
%
%Fixed in flight by enabling {\bf priority inheritance} on that mutex.
%
%\end{frame}

%----------------------------------------
\begin{frame}
  \frametitle{The idea of model checking} \small

\begin{block}{}\small
A model checker takes a {\bf finite-state model} $M$ and a {\bf property} $\varphi$,
and decides whether
\[
M \models \varphi
\]
If not, it returns a {\bf counterexample}: a concrete execution that violates $\varphi$.
\end{block}

\pause
\vspace{0.3cm}

\begin{center}\small
\renewcommand{\arraystretch}{1.25}
\begin{tabular}{|l|p{3.5cm}|p{3.9cm}|}
\hline
\rowcolor{blue!10}
& {\bf Testing} & {\bf Model checking}\\
\hline
{\bf What runs}  & The real system        & A model of the system\\
\hline
{\bf Coverage}   & One execution per test & All executions\\
\hline
{\bf Guarantee}  & Finds bugs, cannot prove absence & Proves absence, within the model\\
\hline
{\bf Scale}      & Full systems           & Limited by the state space\\
\hline
{\bf Effort}     & No model needed        & A model must be built\\
\hline
\end{tabular}
\end{center}

\end{frame}
